% In this file you should put macros to be used only by
% the printed version. Of course they should have a corresponding
% version in macros/web.tex.
% Typically the printed version could have more fancy decorations.
% This should be a very short file.
%
% This file starts with dummy macros that ensure the pdf
% compiler will ignore macros provided by plasTeX that make
% sense only for the web version, such as dependency graph
% macros.


% Dummy macros that make sense only for web version.
\newcommand{\lean}[1]{}
\newcommand{\discussion}[1]{}
\newcommand{\leanok}{}
\newcommand{\mathlibok}{}
\newcommand{\notready}{}
% Make sure that arguments of \uses and \proves are real labels, by using invisible refs:
% latex prints a warning if the label is not defined, but nothing is shown in the pdf file.
% It uses LaTeX3 programming, this is why we use the expl3 package.
\ExplSyntaxOn
\NewDocumentCommand{\uses}{m}
 {\clist_map_inline:nn{#1}{\vphantom{\ref{##1}}}%
  \ignorespaces}
\NewDocumentCommand{\proves}{m}
 {\clist_map_inline:nn{#1}{\vphantom{\ref{##1}}}%
  \ignorespaces}
\ExplSyntaxOff
% Line breaking for long Lean names and paths (pdf only). Lean identifiers such as
% `Physicslib4.Spectral.Stone.isEssentiallySelfAdjoint_generator` are set with \texttt,
% which TeX never breaks, so they overflow the margin. Allow a break after each `.`, `_`
% and `/` inside \texttt, and give TeX extra stretch before it resorts to an overfull line.
\NewCommandCopy{\bpOrigTexttt}{\texttt}
\ExplSyntaxOn
\tl_new:N \l__bp_texttt_tl
\RenewDocumentCommand{\texttt}{m}
 {
  \tl_set:Nn \l__bp_texttt_tl {#1}
  \tl_replace_all:Nnn \l__bp_texttt_tl {\_} {\_\allowbreak}
  \tl_replace_all:Nnn \l__bp_texttt_tl {.} {.\allowbreak}
  \tl_replace_all:Nnn \l__bp_texttt_tl {/} {/\allowbreak}
  \exp_args:NV \bpOrigTexttt \l__bp_texttt_tl
 }
\ExplSyntaxOff
\setlength{\emergencystretch}{3em}
\tolerance=2000
